% ════════════════════════════════════════════════════════════════════
% Chapter 8 — Dataset Format and Metadata Schema
% ════════════════════════════════════════════════════════════════════
\section{Dataset Format and Metadata Schema}\label{sec:dataset}

A research-grade dataset requires not only accurate raw signals but also comprehensive, machine-readable contextual metadata.
Without knowing the athlete's training history, the prescribed load, the subjective fatigue state, and the exact software version that produced the data, no amount of statistical sophistication can produce a reproducible analysis.
This section formalises the Schema~v4 metadata taxonomy and directory layout.


\subsection{Directory Structure}

Each session is stored in a self-contained directory.
Two layout modes are supported:

\paragraph{Flat Mode (Default):}
\begin{lstlisting}[numbers=none, basicstyle=\ttfamily\footnotesize, frame=none, backgroundcolor=\color{white}]
dataset_root/
  session_20260504_143723/
    imu/
      raw_imu.csv
      raw_imu.bin
      camera_imu.csv
    camera/
      marker_positions.csv
      ir_video.mp4
      video_frames.csv
    rotation.json
    smoothed.csv
    annotation_live.csv
    annotation_online.csv
    annotation_offline.csv
    annotation_reviewed.csv
    audit_post_session.png
    metadata.json
    manifest.json
    events.jsonl
    CHECKSUMS.sha256
  ground_truth.csv
\end{lstlisting}

\paragraph{Measurement and derivation are separated by file permission, not by
directory nesting.} Within a session, \texttt{camera/} and \texttt{imu/} hold the
measurement and are read-only with SHA-256 checksums; every other file is derived from
them and can be deleted and rebuilt in a single pass. \texttt{ground\_truth.csv} is the
released label set: one row per accepted repetition, 45 columns, self-contained so that
no downstream consumer needs to open a session directory. Repetitions a reviewer refused
are absent from it, and the refusal itself is recorded in that session's
\texttt{annotation\_reviewed.csv} rather than being a silent deletion.

\paragraph{BIDS Mode (Multi-Subject Archiving):}
When enabled, the directory follows the Brain Imaging Data Structure (BIDS) convention adapted for movement science:
\begin{lstlisting}[numbers=none, basicstyle=\ttfamily\footnotesize, frame=none, backgroundcolor=\color{white}]
dataset_root/
  sub-A7F3/
    ses-2026-05-04/
      run-01_squat/
        imu/  camera/  annotations/  calibration/
        session_info.json
        manifest.json
        events.jsonl
\end{lstlisting}

The subject identifier is the SHA-256 hash of the participant ID, enabling dataset publication with anonymised file paths while maintaining internal consistency for longitudinal studies.


\subsection{Schema Versioning}

Every \file{session\_info.json} and \file{vbt\_config.json} carries a top-level \code{schema\_version} integer.
The current version is~4.
The loading logic enforces the following invariant:
\begin{itemize}[nosep]
  \item Configs with \code{schema\_version > CURRENT\_SCHEMA} are refused with a clear error message, preventing silent misinterpretation of newer-format fields.
  \item Configs with \code{schema\_version <= CURRENT\_SCHEMA} are loaded using \code{\_WITH\_DEFAULT\_} macros, which populate missing fields with struct defaults. This allows sessions recorded under Schema~v2 to be loaded and re-exported by Schema~v4 software without data loss.
\end{itemize}


\subsection{Session Metadata Taxonomy: \code{SessionInfo}}

The \code{SessionInfo} struct contains over 120 fields, organised into nine logical groups.
\Cref{tab:metadata_groups} provides an overview; subsequent subsections detail the non-trivial groups.

\begin{table}[H]
  \centering
  \caption{Overview of \code{SessionInfo} metadata groups.}
  \label{tab:metadata_groups}
  \begin{tabular}{lrl}
    \toprule
    \textbf{Group} & \textbf{Fields} & \textbf{Primary Content} \\
    \midrule
    Identity           & 6   & Session ID, timestamp, schema version, format. \\
    Subject            & 14  & Anonymised hash, sex, age, mass, height, training experience. \\
    Exercise           & 8   & Exercise name, variant, depth criterion, tempo, exercise profile. \\
    Training Context   & 13  & Prescription (\%1RM, target reps/sets), mesocycle phase, RIR. \\
    Conditions         & 6   & Time of day, freshness rating, warmup status, temperature, humidity. \\
    Equipment          & 10  & Bar mass/make, plate manufacturer, rack ID, IMU mount details. \\
    Calibration Prov.  & 8   & Calibration file paths, timestamps, gyro biases, sync residual stats. \\
    Build Provenance   & 8   & Software version, git SHA, branch, dirty flag, build timestamp. \\
    Operator           & 4   & Operator ID, notes, preflight override audit trail. \\
    Sets               & $n\times 11$ & Per-set: weight, target reps, actual reps, timing, RPE, quality. \\
    \bottomrule
  \end{tabular}
\end{table}


\subsection{Subject Metadata: \code{SubjectInfo}}

Subject data is stored in a separate file indexed by the anonymised subject hash, enabling the separation of personally identifiable information (PII) from session data for ethical data management.

\begin{longtable}{llp{7cm}}
  \caption{SubjectInfo fields (stored in \code{sub\_registry/}).} \label{tab:subject_info} \\
  \toprule
  \textbf{Field} & \textbf{Type} & \textbf{Description} \\
  \midrule
  \endfirsthead
  \multicolumn{3}{c}{\textit{Continued from previous page}} \\
  \toprule \textbf{Field} & \textbf{Type} & \textbf{Description} \\ \midrule
  \endhead
  \code{subject\_hash}         & string  & SHA-256 of the participant identifier. \\
  \code{sex}                   & string  & \code{"male"}, \code{"female"}, or \code{"other"}. \\
  \code{age\_years}            & int     & At the time of data collection. \\
  \code{body\_mass\_kg}        & float   & Session-day weigh-in. \\
  \code{height\_cm}            & float   & Measured or self-reported. \\
  \code{training\_exp\_years}  & float   & Total years of resistance training experience. \\
  \code{dominant\_side}        & string  & \code{"left"}, \code{"right"}, or \code{"ambidextrous"}. \\
  \code{arm\_length\_cm}       & float   & Required for press exercises (determines lockout height). \\
  \code{thigh\_length\_cm}     & float   & Required for squat depth estimation from bar position. \\
  \code{tibia\_length\_cm}     & float   & Affects squat kinematics (ankle mobility compensation). \\
  \code{recent\_1rm}           & map\textless string,float\textgreater & Exercise $\rightarrow$ 1RM. E.g., \code{\{"back\_squat": 180.0\}}. \\
  \code{injury\_notes}         & string  & Free-form current injury/limitation text. \\
  \code{consent\_version}      & string  & Version identifier of the signed informed consent form. \\
  \code{ethics\_protocol}      & string  & Institutional ethics protocol number. \\
  \code{data\_sharing\_tier}   & string  & \code{"none"}, \code{"anonymised"}, or \code{"full"}. \\
  \bottomrule
\end{longtable}


\subsection{Training Context: \code{SetInfo}}

Each session may contain one or more working sets.
The \code{SetInfo} struct captures:

\begin{longtable}{llp{6.5cm}}
  \caption{SetInfo fields (per set within a session).} \label{tab:set_info} \\
  \toprule \textbf{Field} & \textbf{Type} & \textbf{Description} \\ \midrule
  \endfirsthead
  \code{set\_id}              & int    & 1-indexed within the session. \\
  \code{weight\_kg}           & float  & Total barbell mass (bar + plates). \\
  \code{target\_reps}         & int    & Programmed repetitions for this set. \\
  \code{target\_rpe}          & float  & Programmed Rating of Perceived Exertion (6--10 Borg CR10). \\
  \code{actual\_rpe}          & float  & Operator-reported RPE after completion. \\
  \code{velocity\_stop}       & string & \code{"none"}, \code{"vl20"} (20\% VL cutoff), \code{"absolute\_0.3"} (cut at $<\SI{0.3}{\metre\per\second}$). \\
  \code{completed\_reps}      & int    & Actual reps detected (from segmenter). \\
  \code{t\_start\_unified\_s} & double & Wall-clock start time. \\
  \code{t\_end\_unified\_s}   & double & Wall-clock end time. \\
  \code{quality}              & string & \code{"good"}, \code{"suspect"}, \code{"failed"}, \code{"discarded"}. \\
  \code{notes}                & string & Free-form operator notes. \\
  \bottomrule
\end{longtable}


\subsection{Equipment Provenance: \code{EquipmentInfo}}

Following the recommendation of Weakley et al.~\cite{weakley2021} that IMU placement significantly affects measurement accuracy, the equipment block documents:

\begin{longtable}{llp{6cm}}
  \caption{EquipmentInfo fields.} \label{tab:equip} \\
  \toprule \textbf{Field} & \textbf{Type} & \textbf{Description} \\ \midrule
  \endfirsthead
  \code{bar\_make}             & string  & Manufacturer and model. \\
  \code{bar\_mass\_kg}         & float   & Nominal mass (e.g., 20.0). \\
  \code{plate\_manufacturer}   & string  & Plate type affects drop dynamics. \\
  \code{rack\_id}              & string  & Station identifier for multi-rack facilities. \\
  \code{imu\_mount\_location}  & string  & \code{"left\_collar"}, \code{"right\_collar"}, \code{"sleeve\_end"}, \code{"center\_knurl"}. \\
  \code{imu\_mount\_orientation} & string & \code{"z\_up"}, \code{"y\_up"}, etc. Determines the gravity alignment axis. \\
  \code{camera\_distance\_m}   & float   & Camera-to-barbell distance. Affects depth noise ($\propto z^2$). \\
  \code{camera\_mount\_type}   & string  & \code{"tripod"}, \code{"wall\_mount"}, \code{"floor"}. \\
  \code{lighting\_condition}   & string  & \code{"natural"}, \code{"artificial"}, \code{"mixed"}. \\
  \code{marker\_type}          & string  & \code{"retroreflective\_tape"}, \code{"retroreflective\_ball"}. \\
  \bottomrule
\end{longtable}


\subsection{Calibration Provenance}

To ensure traceability from kinematic output to calibration state, each session records:

\begin{itemize}[nosep]
  \item Absolute path and timestamp of the accelerometer calibration file (\file{accel\_calib.json}).
  \item Absolute path and timestamp of the camera extrinsic calibration.
  \item The specific gyroscope bias vector used ($[\delta\omega_x,\; \delta\omega_y,\; \delta\omega_z]$, in \si{\degree\per\second}).
  \item Tap-test sync residual: offset (\si{\micro\second}), drift (ppm), and correlation quality.
  \item Placeholder for an Allan-variance characterisation file (for future IEEE~952-2020 compliance).
\end{itemize}


\subsection{Event Log (\code{events.jsonl})}

Every session generates an append-only, newline-delimited JSON event log.
Each event carries a timestamp, event type, and free-form data payload.
The event types include:
\begin{multicols}{2}
\begin{itemize}[nosep]
  \item \code{session\_created}
  \item \code{session\_started}
  \item \code{session\_stopped}
  \item \code{session\_saved}
  \item \code{set\_advanced}
  \item \code{tap\_test\_completed}
  \item \code{sync\_rearm\_triggered}
  \item \code{preflight\_override}
  \item \code{calibration\_completed}
  \item \code{rep\_manual\_mark}
  \item \code{rep\_deleted}
  \item \code{operator\_note}
\end{itemize}
\end{multicols}

The \file{manifest.json} file includes a SHA-256 hash of the entire \file{events.jsonl}, providing tamper detection for the audit trail.


\subsection{Anonymisation Pipeline}

The \file{scripts/anonymise\_session.py} script produces a publication-ready copy of any session by:
\begin{enumerate}[nosep]
  \item Replacing \code{subject\_id} with its SHA-256 hash.
  \item Redacting operator notes, operator ID, and location fields.
  \item Stripping operator-initiated events from \file{events.jsonl}.
  \item Refusing to overwrite an existing output directory (to prevent accidental re-anonymisation).
\end{enumerate}

The anonymised session retains all kinematic data, timing information, and non-PII metadata, making it suitable for deposit on Zenodo or similar open-data repositories.
